\bookchapter{Chapitre 27 : Smartphone facile pour vos parents — guide complet d'accessibilité pour les seniors}

\chapterintro{Honnêtement, en écrivant ce chapitre, j'ai pensé à mes parents. Ma mère, à chaque fois qu'une « mise à jour système » apparaît, n'ose pas appuyer, de peur de casser le smartphone. Mon père avait la police WeChat plus petite que des fourmis, mais il ne savait pas comment l'agrandir, et il a plissé les yeux pendant deux ans. Ce n'est pas qu'ils sont « bêtes », c'est que le smartphone considère par défaut que tout le monde a une vingtaine d'années. Dans ce chapitre, j'utilise les mots les plus simples pour résoudre un par un les problèmes les plus casse-tête des seniors avec leur smartphone.}

\booksection{D'abord, le bon état d'esprit : ce n'est pas « les seniors ne savent pas utiliser », c'est « le smartphone n'est pas bien conçu »}

Avant tout réglage, une idée doit être claire :\textbf{Les interfaces des smartphones actuels n'ont guère pris en compte les utilisateurs seniors dès la conception}. Les pop-ups sont trop petits pour être lus, les gestes nécessitent des touchers précis, les icônes se ressemblent toutes — ce n'est pas un problème d'utilisateur, c'est un problème de produit. Donc, pour aider vos aînés avec le smartphone, la première chose n'est pas d'enseigner, mais de les aider à\textbf{régler l'environnement}. Une fois l'environnement réglé, la moitié des problèmes disparaissent.

Trois règles d'or pour accompagner les seniors avec le smartphone : (1)\textbf{Ne pressez pas}— vous utilisez un smartphone depuis dix ans, eux peut-être depuis trois mois. (2)\textbf{Utilisez des analogies}— « le Wi-Fi c'est comme les tuyaux d'eau à la maison, ouvrez le robinet pour avoir de l'eau », ce genre de comparaison est infiniment plus compréhensible que le jargon technique. (3)\textbf{Écrivez-le}— enseigner oralement une fois peut être oublié, écrire les étapes clés sur un papier collé au dos de la coque du smartphone est plus efficace que tout.

\booksection{I. Première étape : activez le « mode simplifié » — le smartphone devient épuré}

Les smartphones actuels intègrent presque tous un « mode simplifié » (parfois appelé « mode soins », « mode senior », « mode aînés »). Une fois activé, tout le smartphone change du tout au tout :

\begin{itemize}
  \item \textbf{La police et les icônes deviennent très grandes}— fini les lunettes de lecture pour regarder le smartphone.
  \item \textbf{Les icônes du bureau sont grandement simplifiées}— seuls Téléphone, SMS, Appareil photo, Album, WeChat et quelques autres applications courantes sont conservés, le reste est rangé.
  \item \textbf{Les opérations sont devenues plus simples}— de nombreuses fonctions superfétatoires (page -1, assistant intelligent, fenêtres flottantes, etc.) sont désactivées par défaut, l'interface est épurée.
  \item \textbf{Le son est plus fort}— la sonnerie d'appel, le volume des notifications et le son d'appel sont automatiquement réglés au maximum.
\end{itemize}

\booksubsection{Méthodes d'activation du mode simplifié par marque (suivez directement les instructions) :}

\begin{center}
\begin{tabularx}{\linewidth}{|X|X|X|}
\hline
\textbf{Marque} & \textbf{Nom du mode} & \textbf{Méthode d'activation} \tabularnewline
\hline
iPhone & Accessibilité (pas de mode simplifié distinct, mais réglage individuel possible) & Réglages > Affichage et luminosité > Taille du texte, glisser au maximum. Puis « Accessibilité > Affichage et taille du texte » activer « Texte en gras » « Forme des boutons » \tabularnewline
Huawei/Honor & Mode simplifié & Réglages > Système et mises à jour > Mode simplifié > Activer \tabularnewline
Xiaomi / Redmi & Mode minimal & Réglages > Plus de réglages > Mode minimal > Activer \tabularnewline
OPPO / OnePlus & Mode simplifié & Réglages > Fonctions spéciales > Mode simplifié \tabularnewline
vivo/iQOO & Bureau simplifié & Réglages > Bureau, écran de verrouillage et fond d'écran > Bureau simplifié \tabularnewline
Samsung & Mode simplifié & Réglages > Affichage > Mode simplifié \tabularnewline
\hline
\end{tabularx}
\end{center}

Si vous utilisez une autre marque, saisissez « simplifié » ou « senior » dans la barre de recherche en haut des « Réglages » du smartphone, et vous trouverez généralement. Si vous ne trouvez pas, suivez le plan « réglage manuel » ci-dessous, l'effet sera le même.

\booksection{II. Sans mode simplifié, le réglage manuel fonctionne aussi — « grande police + grandes icônes » en cinq étapes}

Certains seniors ne veulent pas du mode simplifié, trouvant l'interface trop « vieille ». Dans ce cas, réglez manuellement les cinq éléments suivants, l'effet est tout aussi bon :

\begin{enumerate}
  \item \textbf{Police au maximum}: Réglages > Affichage > Taille de la police, glisser complètement à droite. Profitez-en pour régler aussi la « taille d'affichage » (parfois appelée « zoom d'affichage ») au maximum, ainsi tous les éléments de l'interface seront agrandis.
  \item \textbf{Icônes agrandies}: Sur l'écran d'accueil, pincez avec deux doigts (comme pour réduire une photo) pour accéder aux réglages du bureau > disposition du bureau (ou taille des icônes), choisissez la disposition avec le moins de cases (par exemple 3×5), les icônes seront naturellement plus grandes.
  \item \textbf{Activer « texte à contraste élevé »}: Rend les bords du texte plus nets, cet interrupteur se trouve dans « Accessibilité ».
  \item \textbf{Désactiver les notifications inutiles}: Réglages > Notifications, désactivez toutes les notifications d'applications sauf Téléphone, SMS et WeChat. Ainsi, le smartphone ne sonnera pas toutes les quelques minutes et ne sera pas agaçant.
  \item \textbf{Placer les contacts fréquents sur le bureau}: Appui long sur un espace vide du bureau > Ajouter un widget > Contacts, placez les appels rapides de votre conjoint, enfants, médecin et contacts d'urgence sur le bureau. Un seul toucher pour appeler, sans parcourir le répertoire.
\end{enumerate}

\begin{tcbpractice}{Exercice : « décorez » le smartphone de vos aînés}
\begin{enumerate}
  \item Activez le mode simplifié (selon le tableau des marques ci-dessus).
  \item Activez l'affichage numérique de la batterie (Réglages > Batterie > Pourcentage de batterie), pour qu'ils sachent toujours combien il reste.
  \item Ouvrez le panneau de notifications, retirez ou masquez les interrupteurs rapides comme Bluetooth, NFC et mode avion dont les seniors n'ont pas besoin.
  \item Placez quatre appels rapides sur le bureau : conjoint, enfants, médecin de famille (si applicable), communauté/gestion.
\end{enumerate}

\end{tcbpractice}

\booksection{III. Prévention de la fraude — lisez ce chapitre avec vos aînés}

Les seniors sont les principales victimes de la fraude téléphonique. Ce n'est pas qu'ils sont moins vigilants, mais les discours des escrocs sont minutieusement conçus pour cibler les faiblesses psychologiques des personnes âgées. Voici les types d'arnaques les plus courants et les parades, que je vous suggère de lire à voix haute à vos aînés ou d'imprimer et de coller près du téléphone.

\booksubsection{Six arnaques courantes (raccrochez immédiatement si vous les repérez) :}

\begin{itemize}
  \item \textbf{(1) « Votre petit-enfant a eu un accident, envoyez vite de l'argent »}— L'escroc dira « votre petit-enfant a fait des bêtises dehors, a été arrêté par la police, il faut une caution ». Retenez :\textbf{Tout appel exigeant un virement pour « libérer quelqu'un » est une arnaque}. La police et les tribunaux ne demandent jamais d'argent par téléphone.
  \item \textbf{(2) « Votre carte d'assurance maladie a été suspendue »}— Le SMS indique que la carte est anormale et vous demande de cliquer sur un lien pour « l'activer ». Retenez :\textbf{L'assurance maladie et la sécurité sociale n'envoient pas de SMS avec des liens à cliquer}. Dans ce cas, appelez directement le 12333 (hotline officielle de la sécurité sociale) pour vérifier.
  \item \textbf{(3) « Placement à haut rendement, sans risque »}— Des groupes WeChat ou SMS recommandent divers « placements de retraite », « projets soutenus par l'État », promettant des rendements élevés. Retenez :\textbf{Tout placement promettant un rendement annuel supérieur à 5 \% avec « garantie en capital » est à 99,9 \% une arnaque}.
  \item \textbf{(4) « C'est ton chef/ton vieil ami, j'ai changé de numéro »}— Se fait passer pour une connaissance pour emprunter de l'argent. Retenez :\textbf{Appeler pour vérifier est plus efficace que tout}, celui qui dit « je ne peux pas prendre d'appel » est à 80 \% un escroc.
  \item \textbf{(5) « Œufs gratuits / conférence sur les compléments alimentaires »}— Des activités gratuites hors ligne attirent les seniors vers une « conférence santé », puis vendent des compléments à prix d'or. Retenez :\textbf{Il n'y a pas de repas gratuit, derrière chaque activité gratuite il y a un business}.
  \item \textbf{(6) « Votre smartphone est infecté, il faut une réparation à distance »}— L'escroc vous fait télécharger une application de contrôle à distance (comme Sunflower, TeamViewer), puis prend le contrôle de votre smartphone pour transférer votre argent. Retenez :\textbf{Ne laissez jamais personne prendre le contrôle à distance de votre smartphone !}Les fabricants de smartphones, les opérateurs de téléphonie et les banques ne vous appelleront jamais pour vous faire installer un logiciel de contrôle à distance.
\end{itemize}

\booksubsection{Trois lignes de défense, une fois en place, elles bloquent 90 \% des fraudes :}

\begin{enumerate}
  \item \textbf{Première ligne : interception intégrée au smartphone}— Sur Android, activez « filtrage des appels indésirables » (dans l'application Téléphone ou Gestionnaire du smartphone) ; sur iPhone, activez « Silencer les appels inconnus » (Réglages > Téléphone > Silencer les appels inconnus).
  \item \textbf{Deuxième ligne : interception par l'opérateur}— Envoyez « KTFSR » par SMS au 10086 pour activer gratuitement le filtrage des appels harcelants fréquents chez China Mobile ; China Telecom et China Unicom proposent des services similaires. Appelez le service client de votre opérateur pour une activation manuelle.
  \item \textbf{Troisième ligne : application nationale anti-fraude}— Téléchargez l'application « Centre national anti-fraude » dans la boutique d'applications, inscrivez-vous et activez l'alerte d'appel. Cette application est reliée à la base de données de la police et peut identifier les numéros d'escrocs dès qu'ils appellent.
\end{enumerate}

\begin{tcbwarn}{« Formule anti-fraude » pour les seniors (imprimez-la et collez-la près du téléphone)}
\begin{itemize}
  \item \textbf{Appel inconnu demandant un virement Raccrochez !}
  \item \textbf{SMS avec lien à cliquer Supprimez !}
  \item \textbf{Promesse de haut rendement garanti Bloquez !}
  \item \textbf{Demande d'installer un logiciel de contrôle à distance Raccrochez !}
  \item \textbf{La vraie police ne fait pas d'enquête par téléphone et ne demande jamais d'argent !}
  \item \textbf{En cas de doute, appelez d'abord vos enfants pour demander.}
\end{itemize}

\end{tcbwarn}

\booksection{IV. Assistance à distance — aidez vos parents avec leur smartphone depuis l'étranger}

Parfois, les problèmes rencontrés par les seniors sont difficiles à résoudre au téléphone car ils ne trouvent pas le bouton correspondant. La fonction « assistance à distance » est alors utile — vous pouvez voir l'écran de leur smartphone directement depuis le vôtre et les aider. Les solutions varient selon les marques :

\begin{itemize}
  \item \textbf{Huawei/Honor}: L'application « MeeTime » intégrée propose une fonction « partage d'écran ». Vous et vos parents installez MeeTime, lancez un appel vidéo puis touchez « Partager l'écran », vous pouvez alors dessiner des cercles sur leur écran pour indiquer à votre mère « touche ici ».
  \item \textbf{Xiaomi / Redmi}: L'application d'appel vidéo Xiaomi propose une « assistance à distance », avec partage d'écran + appel vocal synchronisé.
  \item \textbf{OPPO/vivo}: « Espace familial » ou « garde à distance » intégrés, recherchez « à distance » dans les réglages.
  \item \textbf{Solution universelle (toutes marques)}: Téléchargez l'application « Sunflower Remote Control » (la version gratuite suffit), installez-la chacun de votre côté. Créez un compte commun ou ajoutez-vous en ami, demandez à vos parents de vous dicter le code de vérification affiché sur leur smartphone, et vous pourrez voir et contrôler leur écran à distance.
  \item \textbf{Utilisateurs iPhone}: FaceTime prend en charge le « partage d'écran ». Pendant un appel FaceTime, touchez le bouton de partage d'écran, l'autre personne peut voir votre écran. Mais les restrictions de sécurité d'iOS ne permettent pas le contrôle à distance, seulement un « enseignement dirigé ».
\end{itemize}

\booksection{V. Adaptation de WeChat pour les seniors — quelques phrases suffisent}

WeChat est l'application la plus utilisée par les seniors, mais l'interface est trop dense et complexe. Voici quelques réglages immédiatement efficaces :

\begin{itemize}
  \item \textbf{Agrandir la police de WeChat}: WeChat > Moi > Réglages > Général > Taille de la police, glisser au maximum. Notez que cela n'affecte que la police dans WeChat, pas la police système du smartphone.
  \item \textbf{Activer le « mode soins »}: WeChat > Moi > Réglages > Mode soins > Activer. Après activation, la police et les boutons sont agrandis, le contraste des couleurs est renforcé, et il est possible d'« écouter les messages texte » en un seul toucher (toucher un message texte pour le faire lire à voix haute).
  \item \textbf{Aidez-les à nommer les contacts importants}: Changez le nom des enfants et des parents dans le répertoire pour qu'ils soient faciles à reconnaître — par exemple « Grand fils - Xiaoming », « Deuxième fille - Xiaohong », pour qu'en ouvrant WeChat, ils sachent immédiatement qui est qui.
  \item \textbf{Apprenez-leur à utiliser la « saisie vocale » plutôt que de taper}: Sur le clavier à côté de la zone de discussion, maintenez la barre d'espace (Android) ou touchez l'icône du microphone (iPhone), parlez directement et vos mots se transforment en texte envoyé. Le taux de précision de reconnaissance est désormais très élevé, c'est un confort « qualitatif » pour les seniors qui ne maîtrisent pas le pinyin.
  \item \textbf{Désactiver les fonctions inutilisées dans la « page Découvrir »}: WeChat > Moi > Réglages > Général > Gestion de la page Découvrir, désactivez tout ce qui n'est pas utilisé dans « Moments », « Chaînes vidéo », « Live », « Shopping », « Jeux », pour réduire les distractions.
\end{itemize}

\booksection{VI. Smartphone perdu ? Préparez un « plan d'urgence » pour vos aînés}

Mieux vaut prévenir que guérir. Tant que le smartphone n'est pas perdu, faites d'abord ces quelques choses :

\begin{enumerate}
  \item \textbf{Vérifiez que « Localiser mon appareil » est activé}: iPhone : Réglages > Votre nom > Localiser > Localiser mon iPhone ; Android : recherchez « localiser l'appareil » dans les réglages. Assurez-vous que cet interrupteur est activé.
  \item \textbf{Notez les mots de passe et comptes importants sur papier}: Les mots de passe Apple ID / compte Huawei / compte Xiaomi, une copie enfermée dans un tiroir à la maison, une autre confiée à un enfant de confiance. Après la perte du smartphone, il faut se connecter à ces comptes pour le verrouiller et le localiser à distance ; si vous oubliez le mot de passe, c'est problématique.
  \item \textbf{Apprenez une phrase aux seniors}:\textbf{« Si le smartphone est perdu, la première chose à faire est d'appeler la famille. »}La famille peut immédiatement se connecter à iCloud ou à la page web de localisation d'appareil Android, et le marquer à distance comme « mode perdu », ainsi la personne qui le trouve ne peut ni l'ouvrir ni l'utiliser.
\end{enumerate}

\booksection{VII. Compatibilité écouteurs sans fil et appareils auditifs — une fonction négligée}

Beaucoup de seniors ont des problèmes d'audition, mais ignorent que le smartphone peut fonctionner avec un appareil auditif. L'iPhone et certains smartphones Android prennent en charge les appareils auditifs MFi (Made for iPhone), et peuvent transmettre directement le son de l'appel à l'appareil auditif, bien plus clair que de tenir le combiné du smartphone.

\begin{itemize}
  \item \textbf{Méthode de connexion}: iPhone : Réglages > Accessibilité > Appareils auditifs, activez le Bluetooth, suivez les instructions d'appairage de l'appareil auditif.
  \item \textbf{Sans appareil auditif}: Dans « Accessibilité > Audio/Visuel », activez « Réduction du bruit de l'appel » et « Audio mono » pour améliorer la clarté de l'appel.
  \item \textbf{Android}: Dans « Accessibilité > Amplification auditive », des réglages similaires sont disponibles, l'emplacement varie selon les marques, recherchez directement « appareil auditif ».
\end{itemize}

\booksection{VIII. Fonctions liées à la santé — le smartphone n'est pas qu'un outil de communication}

\begin{itemize}
  \item \textbf{Demande de secours SOS}: Appuyez cinq fois de suite sur le bouton d'alimentation (sur la plupart des smartphones), le smartphone appellera automatiquement les services d'urgence et enverra la localisation aux contacts d'urgence prédéfinis. Aidez les seniors à le configurer à l'avance — iPhone : « Réglages > SOS d'urgence », Android : « Réglages > Sécurité > Secours d'urgence ». Faites une répétition avant de leur apprendre, pour éviter qu'ils n'appuient au mauvais endroit sous le stress.
  \item \textbf{Détection de chute}: L'Apple Watch dispose de cette fonction : si elle détecte une chute suivie d'une période d'immobilité, elle appelle automatiquement les secours et notifie la famille. Si vous en avez le budget, équipez vos aînés, cela peut sauver des vies au moment critique.
  \item \textbf{Rappel de prise de médicaments}: L'application « Santé » de l'iPhone intègre une fonction de rappel de médicaments : définissez le nom du médicament, l'heure et la dose, et un rappel apparaîtra à l'heure voulue. Sur Android, utilisez une « alarme programmée » ou téléchargez une application de type « rappel de médicaments ».
\end{itemize}

\begin{tcbsummary}{Résumé du chapitre}
\begin{itemize}
  \item Activez d'abord le mode simplifié ou agrandissez manuellement la police et les icônes, une fois l'environnement réglé, la plupart des problèmes disparaissent
  \item Trois règles d'or pour accompagner les seniors : ne pas presser, utiliser des analogies, écrire les étapes
  \item Rappelez régulièrement les six arnaques courantes, trois lignes de défense bloquent 90 \% des fraudes
  \item L'assistance à distance vous permet de réparer le smartphone de vos parents depuis l'étranger
  \item Le « mode soins » et la « saisie vocale » de WeChat sont les deux fonctions les plus utiles pour les seniors
  \item Le SOS d'urgence et la détection de chute méritent cinq minutes de configuration, car ils peuvent sauver des vies au moment critique
\end{itemize}

\end{tcbsummary}
